% GENERATED FILE. Edit canonical lessons, then run npm run generate:books.
% canonical-source-hash: fnv1a64:0dfee8ea05b20ed9
% canonical-lessons: TE-C350-neram, TE-C350-nerasthudu, TE-C350-dosi, TE-C350-nirdosi, TE-C350-siksincu

\chapter{Crime and Punishment}
\label{ch:te-crime-and-punishment}
\hlchaptermodality{350}

\begin{chapteropening}
\textbf{By the end of this chapter:} \emph{I can say a crime, a criminal, a guilty person, an innocent person and to punish.}

Complete the last lesson of chapter 350: I can say a crime, a criminal, a guilty person, an innocent person and to punish.
\end{chapteropening}

\section[\texorpdfstring{\te{నేరం} (nēraṁ)}{నేరం (nēraṁ)}]{\te{నేరం} (nēraṁ) — a crime; an offence}

\label{lesson:TE-C350-neram}



\begin{quote}
Before the new one: say the Telugu for a verdict, then the Telugu for bail.
\end{quote}

\subsection*{You'll want to know: \te{నేరం}}
\textbf{\te{నేరం}} — \emph{nēraṁ} — \textquotedblleft{}a crime; an offence\textquotedblright{}.

\begin{quote}
\te{దొంగతనం} \te{పెద్ద} \te{నేరం}.
\end{quote}

\begin{quote}
\emph{doṅgatanaṁ pedda nēraṁ.}
\end{quote}

\begin{quote}
Theft is a serious crime.
\end{quote}

Plural: \te{నేరాలు}, \emph{nērālu}.

\te{నేరం} \te{చేయు}, to commit a crime.

\begin{cousinweb}
\textbf{\te{నేరం}} (\emph{nēraṁ}) is native Dravidian: DEDR 3775 lists Telugu \emph{nēramu}, \textquotedblleft{}fault, offence, crime\textquotedblright{}, with Tamil \emph{nēram}, \textquotedblleft{}guilt, fault\textquotedblright{}.
\end{cousinweb}

\begin{grammarlens}[title={What you've built}]
One more word for telling what happened.
\end{grammarlens}

\subsection*{Your turn}
\begin{itemize}
\raggedright
  \item \emph{Say it:} \emph{nēraṁ}
  \item \emph{Say it:} \emph{nēraṁ}, once more
  \item \emph{Say it:} \emph{doṅgatanaṁ pedda nēraṁ.} — Theft is a serious crime.
  \item \emph{Recall:} say the Telugu for a verdict, then the Telugu for bail, then say \emph{nēraṁ} again
\end{itemize}

\begin{tcolorbox}[breakable,colback=teal!4,colframe=teal!35!black,title={Before you move on}]
What does \te{నేరం} mean? (\textquotedblleft{}A crime; an offence\textquotedblright{}.) Say it once more.
\end{tcolorbox}
\section[\texorpdfstring{\te{నేరస్థుడు} (nērasthuḍu)}{నేరస్థుడు (nērasthuḍu)}]{\te{నేరస్థుడు} (nērasthuḍu) — a criminal; a culprit}

\label{lesson:TE-C350-nerasthudu}



\begin{quote}
Before the new one: say the Telugu for bail, then the Telugu for a crime.
\end{quote}

\subsection*{You'll want to know: \te{నేరస్థుడు}}
\textbf{\te{నేరస్థుడు}} — \emph{nērasthuḍu} — \textquotedblleft{}a criminal; a culprit\textquotedblright{}.

\begin{quote}
\te{పోలీసులు} \te{నేరస్థుడిని} \te{పట్టుకున్నారు}.
\end{quote}

\begin{quote}
\emph{pōlīsulu nērasthuḍini paṭṭukunnāru.}
\end{quote}

\begin{quote}
The police caught the criminal.
\end{quote}

Plural: \te{నేరస్థులు}, \emph{nērasthulu}.

\te{నేరం} with the Sanskrit -\te{స్థ}, standing in: one who stands in a crime.

\begin{cousinweb}
\textbf{\te{నేరస్థుడు}} (\emph{nērasthuḍu}) is a mixed compound: native \emph{nēram}, \textquotedblleft{}crime\textquotedblright{} (DEDR 3775), with the Sanskrit ending \emph{-stha}, \textquotedblleft{}standing in, being in\textquotedblright{} (Monier-Williams).
\end{cousinweb}

\begin{grammarlens}[title={What you've built}]
One more word for telling what happened.
\end{grammarlens}

\subsection*{Your turn}
\begin{itemize}
\raggedright
  \item \emph{Say it:} \emph{nērasthuḍu}
  \item \emph{Say it:} \emph{nērasthuḍu}, once more
  \item \emph{Say it:} \emph{pōlīsulu nērasthuḍini paṭṭukunnāru.} — The police caught the criminal.
  \item \emph{Recall:} say the Telugu for bail, then the Telugu for a crime, then say \emph{nērasthuḍu} again
\end{itemize}

\begin{tcolorbox}[breakable,colback=teal!4,colframe=teal!35!black,title={Before you move on}]
What does \te{నేరస్థుడు} mean? (\textquotedblleft{}A criminal; a culprit\textquotedblright{}.) Say it once more.
\end{tcolorbox}
\section[\texorpdfstring{\te{దోషి} (dōṣi)}{దోషి (dōṣi)}]{\te{దోషి} (dōṣi) — a guilty person; a culprit}

\label{lesson:TE-C350-dosi}



\begin{quote}
Before the new one: say the Telugu for a crime, then the Telugu for a criminal.
\end{quote}

\subsection*{You'll want to know: \te{దోషి}}
\textbf{\te{దోషి}} — \emph{dōṣi} — \textquotedblleft{}a guilty person; a culprit\textquotedblright{}.

\begin{quote}
\te{న్యాయస్థానం} \te{అతన్ని} \te{దోషిగా} \te{తేల్చింది}.
\end{quote}

\begin{quote}
\emph{nyāyasthānaṁ atanni dōṣigā tēlcindi.}
\end{quote}

\begin{quote}
The court found him guilty.
\end{quote}

\te{దోషం} is a fault: \te{దోషి} is the one at fault.

\begin{cousinweb}
\textbf{\te{దోషి}} (\emph{dōṣi}) is borrowed from Sanskrit \emph{doṣin}, \textquotedblleft{}faulty; guilty of an offence\textquotedblright{} (Monier-Williams), from \emph{doṣa}, \textquotedblleft{}a fault\textquotedblright{}.
\end{cousinweb}

\begin{grammarlens}[title={What you've built}]
One more word for telling what happened.
\end{grammarlens}

\subsection*{Your turn}
\begin{itemize}
\raggedright
  \item \emph{Say it:} \emph{dōṣi}
  \item \emph{Say it:} \emph{dōṣi}, once more
  \item \emph{Say it:} \emph{nyāyasthānaṁ atanni dōṣigā tēlcindi.} — The court found him guilty.
  \item \emph{Recall:} say the Telugu for a crime, then the Telugu for a criminal, then say \emph{dōṣi} again
\end{itemize}

\begin{tcolorbox}[breakable,colback=teal!4,colframe=teal!35!black,title={Before you move on}]
What does \te{దోషి} mean? (\textquotedblleft{}A guilty person; a culprit\textquotedblright{}.) Say it once more.
\end{tcolorbox}
\section[\texorpdfstring{\te{నిర్దోషి} (nirdōṣi)}{నిర్దోషి (nirdōṣi)}]{\te{నిర్దోషి} (nirdōṣi) — an innocent person}

\label{lesson:TE-C350-nirdosi}



\begin{quote}
Before the new one: say the Telugu for a criminal, then the Telugu for a guilty person.
\end{quote}

\subsection*{You'll want to know: \te{నిర్దోషి}}
\textbf{\te{నిర్దోషి}} — \emph{nirdōṣi} — \textquotedblleft{}an innocent person\textquotedblright{}.

\begin{quote}
\te{చివరికి} \te{అతను} \te{నిర్దోషి} \te{అని} \te{తేలింది}.
\end{quote}

\begin{quote}
\emph{civariki atanu nirdōṣi ani tēlindi.}
\end{quote}

\begin{quote}
In the end he turned out to be innocent.
\end{quote}

\te{నిర్}-, without, and \te{దోషి}: one without fault.

\begin{cousinweb}
\textbf{\te{నిర్దోషి}} (\emph{nirdōṣi}) is built on Sanskrit \emph{nirdoṣa}, \textquotedblleft{}faultless, guiltless, innocent\textquotedblright{} (Monier-Williams), in the shape of \te{దోషి}.
\end{cousinweb}

\begin{grammarlens}[title={What you've built}]
One more word for telling what happened.
\end{grammarlens}

\subsection*{Your turn}
\begin{itemize}
\raggedright
  \item \emph{Say it:} \emph{nirdōṣi}
  \item \emph{Say it:} \emph{nirdōṣi}, once more
  \item \emph{Say it:} \emph{civariki atanu nirdōṣi ani tēlindi.} — In the end he turned out to be innocent.
  \item \emph{Recall:} say the Telugu for a criminal, then the Telugu for a guilty person, then say \emph{nirdōṣi} again
\end{itemize}

\begin{tcolorbox}[breakable,colback=teal!4,colframe=teal!35!black,title={Before you move on}]
What does \te{నిర్దోషి} mean? (\textquotedblleft{}An innocent person\textquotedblright{}.) Say it once more.
\end{tcolorbox}
\section[\texorpdfstring{\te{శిక్షించు} (śikṣiñcu)}{శిక్షించు (śikṣiñcu)}]{\te{శిక్షించు} (śikṣiñcu) — to punish}

\label{lesson:TE-C350-siksincu}



\begin{quote}
Before the new one: say the Telugu for a guilty person, then the Telugu for an innocent person.
\end{quote}

\subsection*{You'll want to know: \te{శిక్షించు}}
\textbf{\te{శిక్షించు}} — \emph{śikṣiñcu} — \textquotedblleft{}to punish\textquotedblright{}.

\begin{quote}
\te{న్యాయమూర్తి} \te{నేరస్థుడిని} \te{శిక్షించాడు}.
\end{quote}

\begin{quote}
\emph{nyāyamūrti nērasthuḍini śikṣiñcāḍu.}
\end{quote}

\begin{quote}
The judge punished the criminal.
\end{quote}

Past (I …): \te{శిక్షించాను}, \emph{śikṣiñcānu}.

Who is punished takes -\te{ను}; a teacher or a parent can \te{శిక్షించు} too.

\begin{cousinweb}
\textbf{\te{శిక్షించు}} (\emph{śikṣiñcu}) is built on Sanskrit \emph{śikṣā}, \textquotedblleft{}learning, study\textquotedblright{}, which also came to mean \textquotedblleft{}punishment\textquotedblright{} (CDIAL 12432), with the Telugu verb ending \emph{-iñcu}.
\end{cousinweb}

\begin{grammarlens}[title={What you've built}]
One more word for telling what happened.
\end{grammarlens}

\subsection*{Your turn}
\begin{itemize}
\raggedright
  \item \emph{Say it:} \emph{śikṣiñcu}
  \item \emph{Say it:} \emph{śikṣiñcu}, once more
  \item \emph{Say it:} \emph{nyāyamūrti nērasthuḍini śikṣiñcāḍu.} — The judge punished the criminal.
  \item \emph{Recall:} say the Telugu for a guilty person, then the Telugu for an innocent person, then say \emph{śikṣiñcu} again
\end{itemize}

\begin{tcolorbox}[breakable,colback=teal!4,colframe=teal!35!black,title={Before you move on}]
What does \te{శిక్షించు} mean? (\textquotedblleft{}To punish\textquotedblright{}.) Say it once more.
\end{tcolorbox}
